\pagestyle{myheadings}

\chapter{DESCRIPCIÓN DEL PROBLEMA}
En el mundo entero suceden los cambios en los diversos proyectos que se realizan según las normativas propios de cada lugar,~\cite{Aslam2019} menciona que los cambios de diseño son inevitables en los procesos de construcción ocasionando sobrecostes de entre el 5 y el 40 \% del coste del proyecto si se gestionan bien los cambios.

En América Latina los cambios en el alcance tienen un impacto en el costo son muy comunes, por consiguiente la paralización de obras y el aumento presupuestario en América Latina son desafíos significativos que requieren de una gestión eficiente y transparente por parte de los gobiernos. Es fundamental abordar la corrupción, mejorar la planificación y la ejecución de proyectos y buscar soluciones innovadoras para afrontar los desafíos económicos y sociales que enfrenta la región. Además, se debe buscar el equilibrio entre la necesidad de desarrollar infraestructura y el uso responsable de los recursos públicos para garantizar un crecimiento sostenible y equitativo en la región ~\cite[pag.~04]{BID2018}

En el Perú sucede muy a menudo los cambios en el alcance , el cual ocasiona el impacto en el costo a nivel de todo el país.

El contralor General de la República (Nelson Sharck) en ~\cite{CuartoPoder2022} concluye que las causas que ocasionan lo mencionado son la corrupción, la calidad de los expedientes técnicos y una mala gestión de los Proyectos de Inversión Pública por parte de los responsables.

Las consecuencias que trae consigo los cambios en el alcance es cantidad de poblaciones afectadas por obras iniciadas y abandonados aumenta cada día de manera vertiginosa, según ~\cite{CGRP2023} determinando que el 15\% de toda la inversión pública se pierde en sobrecostos, ya sea por problemas de corrupción, problemas de inconducta profesional o ineficiencia, por ende , una mala gestión sumando un total de \$ 29,732,207,833.

El aporte que se dará en el presente trabajo será el análisis de las causas y patrones, evaluación del impacto económico, evaluación del impacto en las partes interesadas y la identificación de las mejores prácticas y lecciones aprendidas.

\section{Problema general}

¿Cual es la relación que existe entre los cambios y los costos en proyectos de inversión pública financiados en el periodo 2018-2022 por el Gobierno Regional de Apurímac?

\section{Problemas específicos}

¿Cual sera la relación que existe entre los cambios basados en el alcance de los trabajos y el  el costo con relación a un objeto de costo en proyectos de inversión publica financiados en el periodo 2018-2022 por el Gobierno Regional de Apurímac?

¿Cual sera la relación que existe entre los cambios basados en el alcance de los trabajos y el  el costo con un patrón de comportamiento en proyectos de inversión pública financiados en el periodo 2018-2022 por el Gobierno Regional de Apurímac?

¿Cual sera la relación que existe entre los cambios basados en el punto de vista de diferentes partes y el  el costo de acuerdo a la relación a un objeto de costo en proyectos de inversión publica financiados en el periodo 2018-2022 por el Gobierno Regional de Apurímac?

¿Cual sera la relación que existe entre los cambios basados en el punto de vista de diferentes partes y el  el costos con un patrón de comportamiento en proyectos de inversión publica financiados en el periodo 2018-2022 por el Gobierno Regional de Apurímac?

¿Cual sera la relación que existe entre los cambios basados en el efecto del cambio y el  el costo de acuerdo a la relación a un objeto de costo en proyectos de inversión publica financiados en el periodo 2018-2022 por el Gobierno Regional de Apurímac?

¿Cual sera la relación que existe entre los cambios basados en el efecto del cambio y el  el costos con un  patrón de comportamiento en proyectos de inversión publica financiados en el periodo 2018-2022 por el Gobierno Regional de Apurímac?


\section{Justificación de la investigación}

\subsection{Justificación teórica}

El estudio propuesto tiene una sólida base teórica al analizar los cambios en el alcance, impacto y costo de proyectos de inversión pública. Se fundamenta en teorías de gestión de proyectos, economía, desarrollo regional y administración pública. Al profundizar en las causas y consecuencias de los cambios en proyectos, se podrá contribuir al conocimiento sobre cómo mejorar la planificación y ejecución de inversiones públicas, identificar factores críticos de éxito y promover buenas prácticas en el ámbito gubernamental.

\subsection{Justificación práctica}

La Sub-Gerencia de Obras del Gobierno Regional de Apurímac ha llevado a cabo una serie de proyectos de inversión pública en el periodo 2018-2022. Estos proyectos son vitales para el desarrollo regional, pero es común que enfrenten cambios en su alcance, lo que puede afectar su impacto y costo. Esta investigación práctica será relevante para la toma de decisiones en futuros proyectos, al proporcionar recomendaciones basadas en el análisis riguroso de experiencias pasadas.

\subsubsection{Justificación Económica}

Los proyectos de inversión pública representan una parte significativa del presupuesto gubernamental. Comprender los factores que influyen en los cambios de alcance e impacto en el costo permitirá una mejor asignación de recursos y una gestión más eficiente de los proyectos. Esta investigación económica será útil para optimizar el uso de los fondos públicos, asegurando que los proyectos alcancen sus objetivos en el marco de las restricciones presupuestarias.

\subsubsection{Justificación Social}
Los proyectos de inversión pública tienen un impacto directo en la calidad de vida de la población. Un análisis en profundidad de los cambios en el alcance e impacto ayudará a comprender cómo estos afectan a las comunidades locales. El estudio permitirá identificar los factores que generan modificaciones en los proyectos y cómo estos cambios afectan a la satisfacción de las necesidades de la población, lo que a su vez puede contribuir a mejorar la participación ciudadana y la transparencia en la gestión pública.

\subsubsection{Justificación Metodológica}
La metodología de investigación propuesta implicará el análisis de datos cuantitativos y cualitativos. Se utilizarán técnicas estadísticas para identificar patrones en los cambios de alcance e impacto en el costo de los proyectos. Asimismo, se recopilarán datos a través de entrevistas, encuestas y análisis de documentos oficiales. La combinación de métodos permitirá obtener una visión holística y enriquecedora de los factores que influyen en los cambios de proyectos de inversión pública.

\subsubsection{Justificación Legal}

El análisis de proyectos de inversión pública también tiene implicaciones legales, relacionadas con la contratación pública, licitaciones, permisos y normativas aplicables. Al comprender cómo los cambios afectan la legalidad de los proyectos, se podrán proponer recomendaciones para fortalecer los procesos legales y garantizar el cumplimiento de la normativa vigente.

\subsubsection{Justificación Investigativa}
Este estudio representa una contribución original al campo de la gestión de proyectos de inversión pública en el contexto regional del Gobierno de Apurímac. No existen investigaciones similares que aborden este tema específico, lo que lo convierte en una valiosa oportunidad para expandir el conocimiento en esta área y sentar las bases para futuras investigaciones y mejoras en la ejecución de proyectos públicos.

En conclusión, el estudio propuesto sobre los cambios en el alcance, el impacto y el costo de los proyectos de inversión pública ejecutados por el Gobierno Regional de Apurímac durante el período 2018-2022 es valioso tanto en términos teóricos como prácticos. Proporcionará conocimientos relevantes para mejorar la gestión pública, optimizar recursos y tomar decisiones más informadas, buscando así el desarrollo sostenible y el beneficio de la comunidad regional.
\section{Ubicación y contextualización}
\subsection{Ubicación}
\subsubsection{Geográfica}
El tema de la tesis se centrará en los proyectos de inversión pública ejecutados por el Gobierno Regional de Apurímac, por lo tanto, la ubicación geográfica para esta investigación se encuentra en la región de Apurímac, Perú. Apurímac es una región situada en la zona sur-central del país, limitando al norte con Ayacucho, al este con Cusco, al sur con Arequipa y al oeste con la región de Ayacucho y el departamento de Cusco.

\subsubsection{Ubicación Temporal}
El período de tiempo que abarcará la investigación es el comprendido entre el año 2018 y el año 2022. Durante estos cinco años, el Gobierno Regional de Apurímac llevó a cabo diversos proyectos de inversión pública con el objetivo de mejorar la infraestructura, servicios y desarrollo en la región. Al enfocarse en este periodo, la tesis analizará un lapso reciente y relevante para comprender los cambios en el alcance, impacto y costo de los proyectos ejecutados en ese periodo específico.

\subsection{Alcances}
\begin{itemize}
	\item Análisis de proyectos de inversión pública: El estudio se enfocará específicamente en proyectos financiados con inversión pública ejecutados por el Gobierno Regional de Apurímac durante el periodo 2018-2022. Se examinarán proyectos de diferentes sectores, como infraestructura vial, educación, salud, agua y saneamiento, entre otros.
	\item Evaluación de cambios en el alcance: Se analizarán los cambios que ocurrieron en el alcance original de los proyectos de inversión pública durante su ejecución. Se buscará identificar las causas y motivaciones detrás de estos cambios, así como sus implicancias en términos de objetivos y resultados esperados.
	\item Impacto en el costo: El estudio indagará sobre el impacto financiero que los cambios en el alcance de los proyectos generaron en su costo total. Se evaluará si estos cambios llevaron a un aumento o disminución de los recursos asignados inicialmente, y se analizará cómo esto afectó la viabilidad y eficiencia de la inversión pública.
	\item Implicancias sociales y económicas: Se considerarán las implicancias sociales y económicas de los cambios en el alcance e impacto en el costo de los proyectos. Se analizará cómo estos cambios afectaron a las comunidades locales, la generación de empleo, el bienestar de la población y el desarrollo regional en general.
\end{itemize}
\subsection{Delimitaciones}
\begin{itemize}
	\item Exclusión de proyectos privados: El estudio se enfocará únicamente en proyectos financiados con inversión pública ejecutados por el Gobierno Regional de Apurímac. Se excluyen proyectos con financiamiento privados o ejecutados por otras entidades gubernamentales u \gls{ong}.

	\item Limitación temporal: El análisis se circunscribe al periodo comprendido entre los años 2018 y 2022. Se excluyen \gls{pip} ejecutados fuera de este intervalo temporal.
	\item Restricción de sectores: Aunque se incluirán proyectos de diferentes sectores, se establecerá un número determinado de sectores para su análisis. Esto implica que algunos sectores podrían quedar excluidos debido a la limitación de recursos y tiempo para la investigación.
	\item Disponibilidad de datos: La investigación estará sujeta a la disponibilidad y accesibilidad de datos relacionados con los proyectos de inversión pública en Apurímac. Si la información necesaria no está disponible o es insuficiente, algunos aspectos del análisis podrían verse afectados.
	\item Alcance regional: La investigación se centrará en los proyectos de inversión pública ejecutados específicamente en la región de Apurímac. No se incluirán proyectos de otras regiones del Perú en el análisis.

\end{itemize}
Con estas delimitaciones y alcances, la investigación sobre los cambios en el alcance e impacto en el costo de proyectos de inversión pública ejecutados por el Gobierno Regional de Apurímac durante el periodo 2018-2022 se llevará a cabo de manera enfocada y rigurosa, permitiendo obtener resultados pertinentes y relevantes para la toma de decisiones y el desarrollo regional.

\subsection{Limitaciones}

\begin{itemize}
	\item Disponibilidad limitada de fuentes: El tema de investigación específico sobre "Cambios en el alcance e impacto en el costo en proyectos de inversión pública ejecutados en el periodo 2018-2022 por el Gobierno Regional de Apurímac" puede tener una cantidad limitada de fuentes bibliográficas disponibles, ya que podría ser un tema específico y de interés localizado.
	\item Actualidad de la información: Dado que el periodo de investigación abarca hasta el año 2022, la información más actualizada disponible podría estar limitada, lo que podría dificultar el acceso a datos más recientes sobre los proyectos ejecutados en ese periodo.
	\item Acceso restringido a informes oficiales: Los datos y análisis detallados sobre los proyectos de inversión pública ejecutados por el Gobierno Regional de Apurímac podrían estar sujetos a restricciones de acceso debido a cuestiones de confidencialidad o políticas gubernamentales.
	\item Idioma: Es posible que algunas fuentes relevantes estén disponibles solo en idiomas locales o en español, lo que podría limitar la accesibilidad para investigadores que no dominen el idioma.
	\item Falta de estudios específicos: Dado que el período de investigación es relativamente reciente y puede estar relacionado con proyectos específicos del Gobierno Regional de Apurímac, podría haber una falta de estudios o investigaciones anteriores que aborden exactamente el mismo tema.
	\item Información incompleta o no actualizada: La información disponible sobre los proyectos de inversión pública ejecutados durante el período 2018-2022 podría ser parcial o no estar actualizada, lo que podría afectar la precisión y completitud de los datos utilizados en la investigación.
\end{itemize}

\subsection{Ética de la investigación}
\begin{itemize}
	\item Consentimiento informado: Los investigadores deben obtener el consentimiento informado de todas las partes involucradas en la investigación. Esto incluye a los funcionarios del Gobierno Regional de Apurímac, así como a cualquier otra persona que pueda verse afectada directa o indirectamente por los resultados del estudio.
	\item Confidencialidad y privacidad: Se debe garantizar la confidencialidad de los datos y la privacidad de las personas involucradas en el estudio. Los datos recopilados deben ser manejados de manera segura y solo utilizados con fines investigativos específicos, protegiendo la identidad y la información personal de los participantes.
	\item Imparcialidad y objetividad: Los investigadores deben mantener la imparcialidad y objetividad en todas las etapas del estudio. Los resultados deben basarse en el análisis de datos precisos y fiables, evitando sesgos y asegurando una interpretación justa y equitativa de los hallazgos.
	\item Integridad académica: La investigación debe conducirse con honestidad y rigor académico. Los investigadores deben evitar la fabricación, falsificación o manipulación de datos, así como el plagio de trabajos previos.
	\item Beneficio y riesgo: Los investigadores deben considerar el equilibrio entre los beneficios potenciales y los riesgos asociados con el estudio. Es fundamental que los posibles beneficios de la investigación superen cualquier daño potencial a los participantes o a la comunidad en general.

	\item Aprobación ética: Si la investigación involucra seres humanos o datos sensibles, los investigadores deben obtener la aprobación de un comité de ética o cualquier entidad responsable de supervisar la integridad y el bienestar de los participantes.

	\item Transparencia y divulgación: Los investigadores deben ser transparentes en la divulgación de sus métodos, resultados y conclusiones. Esto incluye informar sobre cualquier conflicto de intereses y compartir los datos subyacentes para que otros puedan validar y replicar el estudio.
	\item Uso responsable de los resultados: Los resultados de la investigación deben utilizarse de manera responsable y ética. Los investigadores deben considerar las implicaciones de sus hallazgos y su posible impacto en la toma de decisiones y políticas públicas.

	\item Respeto a la cultura y comunidad local: Si la investigación se realiza en una comunidad específica, los investigadores deben respetar y valorar la cultura, tradiciones y conocimientos locales, evitando la apropiación indebida de la información recopilada.

	\item Beneficio social: En la medida de lo posible, la investigación debe buscar contribuir al bienestar social y al desarrollo de la comunidad, proporcionando información relevante y útil para mejorar la gestión de Proyecto de inversión pública en el Gobierno Regional de Apurímac.

\end{itemize}
